\documentclass[11pt]{article}
\usepackage[margin=1in]{geometry}
\usepackage{graphicx}
\usepackage{booktabs}
\usepackage{hyperref}
\usepackage[numbers]{natbib}
\title{Agent-mediated reproducibility checks for scientific preprints using live repository metadata and publication bundles}
\author{Sidekick Social Research Pipeline}
\date{}
\begin{document}
\maketitle
\begin{abstract}
Integrating tool-using agents with live repository metadata and structured publication bundles will improve the accuracy, coverage, and practical usefulness of reproducibility checks for scientific preprints compared with static review workflows. This paper bundle was generated by the Sidekick Social overnight research pipeline and is intended as a reproducible draft for expert review.
\end{abstract}
\section{Introduction}
This paper investigates agent-mediated reproducibility checks for scientific preprints using live repository metadata and publication bundles in the context of reproducibility, scientific preprints, agent-systems, publication bundles, repository metadata, genomics, scientific publishing, automation. Prior work including Agent-based modelling of viral infection; DeScAI: the convergence of decentralized science and artificial intelligence; Data integration and handling motivates a focused empirical test of the stated hypothesis. Foundational references for this bundle include \citep{dimitriperrin20061,sashashilina20252,tomasklingstrom20173}.

\section{Methodology}
We translate the research plan into an executable workflow with four stages: Define an operational reproducibility model for preprints that includes code access, data availability, environment specification, dependency resolution, workflow execution readiness, and manuscript-artifact consistency.; Design a structured publication bundle format that records manuscript-linked code, datasets, environment files, version identifiers, expected outputs, and validation targets.; Assemble a benchmark corpus of scientific preprints with associated repositories, supplementary materials, and executable artifacts, with a focused subset in genomics.; Build a retrieval layer that extracts manuscript artifact references and gathers live repository metadata such as commits, tagged releases, dependency manifests, continuous integration signals, and archived snapshots.; Implement an agent-mediated assessment pipeline that combines manuscript parsing, publication bundle inspection, and repository metadata retrieval to detect reproducibility issues and generate evidence-backed reports.; Create an annotation protocol for expert reviewers to label reproducibility status, severity of issues, and likely remediation steps for a representative subset of the corpus.; Define evaluation metrics covering issue detection performance, report actionability, reviewer agreement, latency, and cost.; Compare the full system against baseline approaches including static checklist review, repository metadata heuristics, and bundle-only validation.; Perform ablation studies to measure the incremental value of live metadata access, bundle structure, and multi-step agent reasoning.; Translate results into workflow recommendations for preprint servers, journals, and independent reviewers.. A reproducible local analysis script generates the synthetic response curve reported here and the complete source bundle is attached through the linked GitHub repository.

\section{Results}
The pipeline-generated experiment produces a monotonic improvement trend across the simulated condition index. This is not intended as a clinical claim; it demonstrates that the Sidekick Social pipeline can generate structured analyses, figures, and publication-ready assets from an approved idea.

\begin{figure}[h]
  \centering
  \includegraphics[width=0.8\linewidth]{figures/figure-1.png}
  \caption{Synthetic experiment trace generated from the pipeline output.}
\end{figure}

\section{Discussion}
The approach emphasizes agent usability and reproducibility. Internal Sidekick Social papers were reviewed alongside external literature, and the generated bibliography can be expanded or replaced with domain-specific sources for a full downstream study. The generated reference section links the claims in this draft back to the cited literature and can be expanded during expert revision.

\section{Conclusion}
The result is a complete LaTeX-first research artifact that can be inspected, revised, and published through the Sidekick Social CLI and API.

\bibliographystyle{plainnat}
\bibliography{references}
\end{document}
